\documentclass[a4paper]{report}

\usepackage[a4paper, total={7in, 10in}]{geometry}
\usepackage[english]{babel}
\usepackage[utf8]{inputenc}
\usepackage[export]{adjustbox}
\usepackage{amsmath}
\usepackage{graphicx}
\usepackage{wrapfig}
\usepackage{subcaption}
\usepackage{epigraph}
% pour la modification du style des pages (numérotation)
%\usepackage{fancyhdr}



\graphicspath{ {images/} }

\title{\LARGE{RAPPORT DE PROJET}}
\author{\textsc{Dressayre} Nicolas, \textsc{Sylla} Cheick}
\date{\today}
\makeatletter

%\pagestyle{fancy}
%\fancyhf{}

%\rhead{RAPPORT DE STAGE__________}
%\lhead{Présentation}
%\renewcommand{\headrulewidth}{0pt}
%\rfoot{Page \thepage \hspace{1pt} sur n}


\begin{document}
	
	
	% --------------------------------------------------- PAGE DE GARDE
		\enlargethispage{2cm}
		
		\begin{center}
			%logos
			\includegraphics[width=6cm]{logo_um.png}
			
			%titre
			\vspace*{3cm}
			\Huge{\@title}\\
			
			
			\large{\textbf{Faculté des sciences de Montpellier}}\\
			%\vspace*{0.5cm}
			Département Électronique, Énergie électrique et Automatique\\
			du 17 Janvier au 30 Mai 2022 \\
			
			%formation
			\vspace*{1cm}
			\large{Présenté pour obtenir la}\\
			\vspace*{0.5cm}
			\large{\textbf{Licence Générale}}\\
			\large{\textbf{Électronique, Énergie électrique et Automatique}}\\
			%\vspace*{0.5cm}
			%\large{A l'Institut Universitaire de Technologie de Tarbes}\\
			
			%auteur
			\vspace*{1cm}
			\large{Par}\\
			\vspace*{0.5cm}
			\large{\@author} 
			
			\vspace*{1cm}
			\large{Devant}\\
			
			\vspace*{1cm}
			\large{\textbf{Enseignent Référent}}\\
			\large{Aurore Vicet}\\
			Maître de conférences - HDR / Associate professor\\
			
			\vspace*{1cm}
			\large{\textbf{Membre du jury ...}}\\
			\large{m. Dupont}\\
			Lorem ipsum dolor sit amet\\
			
			%	\begin{tabular}{lr}
			%	\textbf{Maître de stage dans l'entreprise}					& \textbf{Enseignant Référent} 	\\
			%	Nathalie Yerle 											& Fabien Lacressonniere 			\\
			%	  Traction Control / Software Skill Leader 					& Maître de Conférences à l'IUT de Tarbes	\\
			%	\end{tabular}
					
		\end{center}
	\pagebreak % PAGE BREAK !!! ------------------------------------------------- PAGE BREAK !!!		
	% --------------------------------------------------- FIN DE PAGE DE GARDE
	
	% --------------------------------------------------- DÉBUT ABSTRACT
		
		\begin{center}
			\textbf{Introduction}
		\end{center}
		{\huge C}e document rapporte mon expérience et mes travaux au cours du projet de réalisé en troisième énnée de licence générale en Électronique, Énergie électrique et Automatique réalisée à la Faculté des Sciences de Montpellier.\\
		
		L'objectif premier de ce projet est de ...

		
		\vspace*{3cm}
		
		\begin{center}
			\textbf{Remerciements}
		\end{center}
	
		{\huge N}ous tenons à remercier toutes les personnes qui nous ont apporté leur aide lors ce projet et pendant la rédaction de ce rapport.\\

En premier lieu, nous voudrons exprimer mes remerciements à ...
Nous tenons aussi à remercier ...\\

Nous saisissons cette occasion pour remercier les enseignants du département EEA, notamment.... Tous ont garanti non seulement le bon déroulement du projet, mais aussi de l'année universitaire.\\

%Je remercie également l'équipe Contrôle Commande, en particulier Didier Colin et Alexandre Lamoulie qui ont su répondre à mes questions et procurer les informations nécessaires lorsque j'ai rencontré des difficultés.\\

Pour finir, un grand merci à... %/mes amis et collègues qui m'ont guidé pour la rédaction de ce rapport.
		
%		\vspace*{3cm}
%		\begin{center}
%			\textit{"Les vieux deviennent cons. \\ Les jeunes deviennent vieux".}\\
%			Alexandre L. , 2021
%		\end{center}

	% --------------------------------------------------- FIN ABSTRACT
	
	% --------------------------------------------------- DÉBUT REMERCIEMENTS 
	%remerciements sur meme page que abstract
	% --------------------------------------------------- FIN REMERCIEMENTS
	
	% --------------------------------------------------- DÉBUT SOMMAIRE
	\renewcommand*\contentsname{Sommaire} % changement du titre de la tableofcontents
	\tableofcontents
	\pagebreak
	% --------------------------------------------------- FIN SOMMAIRE
	
	% --------------------------------------------------- DÉBUT LISTE TAB/FIG
	% --------------------------------------------------- FIN LISTE TAB/FIG
	
	% --------------------------------------------------- DÉBUT GLOSSAIRE
	% --------------------------------------------------- FIN LISTE GLOSSAIRE
	
	
	% --------------------------------------------------- DÉBUT DÉVELOPPEMENT
	
	\renewcommand{\chaptername}{Partie} % changement du titre des chapitres en "parties"
	\chapter{Présentation de la spectroscopie et de la QEPAS}
	
	\section{Spectroscopie}
	\section{QEPAS}
	\section{Mise en oeuvre a l'IES}

	\chapter{Contexte et objectifs du projet}

	
	\section{Contexte}
	Le besoin / ce que l'on doit accomplir avec ce projet
	
	\section{Sujet: Afficheur de mesures de concentration de gaz}
	Conception et réalisation du système autours de la Raspi
		
	\section{Sujet secondaire: Comparaison des performances entre QEPAS et capteurs du commerce}
	Mise en application du système, utilisation avec les 2 types de mesures et conclusion
	
	\chapter{Afficheur de mesures de concentration de gaz}
	
	
	\section{Besoin technique}
	
	\section{Gestion de projet et organisation du travail}
	
	\section{Conception}
	\subsection{architecture retenue}
	\subsection{choix des composants}
	\subsection{éléments logiciels}
	\subsection{modélisation du boitier}
	
	\section{Réalisation}
	\subsection{prototypage sur breadboard}
	\subsection{tests préliminaires}
	\subsection{impression du boitier et assemblage}
	
	\section{conclusion, difficultés rencontrées}

	\chapter{Comparaison des performances entre QEPAS et capteurs du commerce}
	\section{Étalonnage des capteurs MQx}
	\subsection{méthode employée}
	\subsection{déroulement}
	\subsection{conclusion (comparaison avec datasheet ?)}
	
	\section{Mesures avec les deux techniques}
	\subsection{méthode employée}
	\subsection{déroulement}
	\subsection{conclusion}
	

	% --------------------------------------------------- FIN DÉVELOPPEMENT	
	\chapter{Valorisation du projet par rapport a la poursuite d'études}
	% --------------------------------------------------- DÉBUT CONCLUSION
	
		Ce projet nous a permit de réaliser... \\
		Nous avons pu mettre en œuvre des connaissances acquises en L3 EEA.\\
	
	%%	Pour finir, nous tenons à dire que nous conservons une très bonne impression de cette courte expérience dans le domaine du ferroviaire.\\
 
	
	% --------------------------------------------------- FIN CONCLUSION
	\renewcommand{\bibname}{Références} % changement du titre des chapitres en "parties"
	\begin{thebibliography}{9}
		
		
		\bibitem{transfer_thermique_pspice} 
		Shuhui Li, Rajab Challoo et Robert A. McLauchlan. 
		\textit{Heat Transfer Simulation Using PSpice}.\\ 
		Conference: ASME 2003 Heat Transfer Summer Conference.\\
		DOI:10.1115/HT2003-47021
		
		\bibitem{brake_resistor} 
		Thomas B. D. Björklund. 
		\textit{A brake resistor perspective on Thermal Management}.\\ 
		CDanotherm Electric [En ligne]. 2015 [consulté le 26 mars 2021].\\
		Disponible sur: \\\texttt{https://www.danotherm.com/power-resistors/downloads}
		
		\bibitem{histoire_alstom} 
		\textit{Génèse et Histoire du site Alstom à Séméac}.\\
		Maire de Soues. [En ligne]. 2014 [consulté le 20 mai 2021].\\
		Disponible sur: \\\texttt{https://www.soues.com/Fichiers/pages/155333145124genese-et-histoire-du-site-alstom-a-semeac.pdf}
		
		\bibitem{ACTGV_pmcf} 
		\textit{PMCF Onduleurs}.\\
		Association des Conducteurs de Trains à Grande Vitesse. [En ligne] M.Durochat et G.Desplanques [consulté le 05 juin 2021].\\
		Disponible sur: \\\texttt{http://actgv.fr/wp-content/uploads/2012/05/PMCF-onduleur.pdf}
		
		\bibitem{ACTGV_puissance} 
		\textit{Circuits de Puissance Traction}.\\
		Association des Conducteurs de Trains à Grande Vitesse. [En ligne] M Durochat, Jean Willemin et G. Desplanques, 2016 [consulté le 03 juin 2021].\\
		Disponible sur: \\\texttt{http://actgv.fr/wp-content/uploads/2016/10/Circuit-de-Puissance-21-04-2016-dernier-travail.pdf}
		
		\bibitem{alstom_site} 
		\textit{Site officiel d'Alstom}.\\
		Site officiel d'Alstom. [En ligne] [consulté le 08 juin 2021].\\
		Disponible sur: \\\texttt{https://www.alstom.com}
		
	\end{thebibliography}
	
	
\end{document}